% Generated draft: requires mathematical and source review.
\documentclass[11pt]{article}
\usepackage[margin=25mm]{geometry}
\usepackage{fontspec}
\setmainfont{Latin Modern Roman}
\usepackage{amsmath,amssymb,mathtools,unicode-math}
\setmathfont{Latin Modern Math}
\usepackage{enumitem,needspace,xurl,hyperref}
\setlength{\parindent}{0pt}
\setlength{\parskip}{5pt}
\setlength{\emergencystretch}{3em}
\newcommand{\N}{\mathbb{N}}
\newcommand{\Z}{\mathbb{Z}}
\newcommand{\Q}{\mathbb{Q}}
\newcommand{\R}{\mathbb{R}}
\newcommand{\C}{\mathbb{C}}
\newcommand{\F}{\mathbb{F}}
\newcommand{\A}{\mathbb{A}}
\newcommand{\PP}{\mathbb{P}}
\newcommand{\E}{\mathbb{E}}
\newcommand{\eps}{\varepsilon}
\newcommand{\dd}{\,\mathrm{d}}
\begin{document}
\section*{UnsolvedMath ID: 30006524}
\textbf{Source problem code:} catalogue entry
\subsection{Definitions and mathematical statement}
In standard geometric terminology, a \(k\)-simplex is the convex hull of \(k+1\) affinely independent points in a real affine space; a triangle is a \(2\)-simplex. These definitions provide terminology only: the supplied record does not establish the intended geometric configuration or counting convention.

The retained source fragment asks: ``Left: How many triangles? Right: How many simplices?'' It refers to Figure~1 in a contribution entitled \emph{Permutahedron Triangulations via Total Linear Stability and the Dual Braid Group}. The figure and its defining parameters are absent from the supplied record. Consequently, no particular set of triangles or simplices, simplex dimension, triangulation, or counting rule is specified, and no standalone mathematical counting problem can be recovered. In particular, the title alone does not specify a permutahedron or a triangulation to count.

The supplied record classifies the statement as \texttt{unrecoverable}. Its \texttt{open} status is a supplied registry label, not a verified assertion that a definite mathematical problem remains unsolved.
\subsection{Short English statement}
The source asks how many triangles and simplices occur in an unavailable figure. Without that figure, its parameters, and a counting convention, the objects to be counted cannot be identified. This fragment is retained for provenance and requires source recovery and human mathematical review before a precise problem can be stated.
\subsection{Sources}
\begin{itemize}
\item Supplied local UnsolvedMath registry record, permanent ID \(30006524\), title \emph{Triangle-Counting Figure Caption}, source problem code \texttt{catalogue entry}, in the Oberwolfach Reports Open Problems collection.
\item Supplied report citation: \emph{Enumerative Combinatorics} (2026), Oberwolfach Reports, pp.~81--156, DOI \texttt{10.4171/owr/2026/2}. \url{https://doi.org/10.4171/owr/2026/2}. The fragment identifies Nathan Williams's contribution \emph{Permutahedron Triangulations via Total Linear Stability and the Dual Braid Group}, described as joint work with Colin Defant and Melissa Sherman-Bennett. No separate problem proposer is supplied.
\item Supplied literature reference: Defant--Sherman-Bennett--Williams, \emph{Permutahedron Triangulations}. \url{https://arxiv.org/abs/2509.11497}.
\item Literature assessment dated exactly \texttt{2026-08-22}: supplied dated curation, not independently verified. It labels the record open while explaining that the exact record is a figure-caption fragment and does not identify a definite claim that can be classified as solved. The supplied background requests independent verification of time-sensitive status and cited claims.
\item The supplied statement-verification field reports that the official report confirms the figure-caption origin and that the missing diagram and parameters prevent recovery of a standalone counting problem. This verification is reported from the supplied record; no source was independently visited or checked for this draft.
\end{itemize}
\end{document}
